\documentclass[11pt]{article}
\usepackage[T1]{fontenc}
\usepackage{lmodern,amsmath,amssymb,amsthm,geometry,microtype,xcolor}
\geometry{margin=0.85in}
\usepackage[colorlinks=true,allcolors=blue!45!black]{hyperref}
\setlength{\emergencystretch}{3em}
\newtheorem{theorem}{Theorem}
\title{A sign correction to a printed\newline nonreducible-operation characterization}
\author{Myeongjin Kang}
\date{October 26, 2026}
\begin{document}\maketitle
\begin{abstract}
Conjecture 5.3 of Kiss and Somlai, recorded as MathDB 336624, prints an
alternating representation with coefficients $(-1)^i$ starting at $i=1$.
For odd arity its diagonal is incompatible with idempotence on a nontrivial
ordered group. We verify the concrete ternary heap $F(x,y,z)=x-y+z$ on the
real line: it is associative, idempotent, strictly monotone in each coordinate,
and cannot be obtained by iterating any associative binary operation.
Nevertheless it has no representation with the printed signs, even allowing
arbitrary bijections. All these assertions, including the source's one-based
finite sum, are proved in Lean. This is a literal sign correction; the
sign-corrected classification is outside its scope.
\end{abstract}
\noindent\textit{Scope and credit.} The heap example is already given in
Example 4.10 of the source. No new example, new classification theorem, or
worldwide priority is claimed. This AI-assisted paper has original prose
licensed CC BY 4.0 and original code licensed MIT.

\section{The printed statement}
Let $X$ be a totally ordered Abelian group. Conjecture 5.3 in
arXiv:1707.04341v3, printed p.~11, equation (24), asserts a characterization
of nonreducible associative idempotent strictly monotone functions of odd
arity by
\begin{equation}\label{eq:printed}
 F(x_1,\ldots,x_n)=g^{-1}\left(\sum_{i=1}^n(-1)^i g(x_i)\right),
\end{equation}
where $g$ is a monotone bijection~\cite{KS}. The same signs appear in the
original PDF and in v1. MathDB 336624 records this formula but omits the
word ``strictly'' from the monotonicity assumption~\cite{MathDB}; the
example below satisfies the stronger v3 assumption as well.

Monotonicity here means that each coordinate map is order-preserving or
order-reversing; different coordinates may have different directions.
For a ternary operation, associativity means that the three placements
of an inner ternary operation in five arguments agree. Reducibility means
that some associative binary operation $G$ satisfies
$F(x,y,z)=G(G(x,y),z)$ for all inputs. No monotonicity or continuity
assumption is placed on the proposed binary reduction.

\section{A complete witness}
\begin{theorem}
On $\mathbb R$ with its usual order and addition, put $F(x,y,z)=x-y+z$.
The operation is ternary associative and idempotent, strictly increasing
in its first and third coordinates, and strictly decreasing in its second.
It has no associative binary reduction. There is no bijection
$g:\mathbb R\to\mathbb R$ satisfying the printed formula
\[
 F(x,y,z)=g^{-1}(-g(x)+g(y)-g(z)).
\]
Thus the literal characterization~\eqref{eq:printed} fails at odd arity three.
\end{theorem}
\begin{proof}
Every nesting of $F$ in five arguments gives $a-b+c-d+e$, proving
associativity. Also $F(x,x,x)=x$, and the strict coordinate monotonicity
follows immediately from the coefficients $1,-1,1$.

Suppose an associative binary operation $G$ were a reduction, and let
$c=G(0,0)$. The reduction identity at $(0,0,z)$ gives $G(c,z)=z$.
Associativity and the identity at $(x,0,0)$ then give
\[
 G(x,c)=G(x,G(0,0))=G(G(x,0),0)=x.
\]
Applying the reduction identity at $(c,x,c)$ therefore yields
\[
 x=G(G(c,x),c)=F(c,x,c)=2c-x\qquad(x\in\mathbb R).
\]
Taking $x=0$ forces $c=0$, while $x=1$ forces $c=1$, a contradiction.
Thus no associative binary reduction exists.

Finally, suppose a bijection $g$ gave the printed representation. On the
diagonal, idempotence and injectivity give
\[
 x=g^{-1}(-g(x)),\qquad g(x)=-g(x),\qquad g(x)=0
 \quad(x\in\mathbb R).
\]
Then $g(0)=g(1)$ contradicts injectivity. This obstruction does not even
require monotonicity of $g$.
\end{proof}

\section{Scope and formal certificate}
The obstruction concerns exactly the signs printed in the source. If
$(-1)^i$ is replaced by $(-1)^{i+1}$, the identity bijection represents the
heap as $x-y+z$; the example then supplies no counterexample. We make no
claim about that intended classification or its status.

The Lean module defines ternary associativity with both nesting equalities,
strict coordinate monotonicity, actual reducibility by an arbitrary binary
operation, and the printed representation by a monotone or antitone
bijection. Its finite sum is indexed by \texttt{Fin 3}, with exponent
\texttt{i.val+1}; a separate lemma proves that it equals
$-g(x_1)+g(x_2)-g(x_3)$. Hence the sign obstruction is tied to the actual
one-based source formula. The declaration \texttt{heap\_not\_reducible}
proves the full quantification over all associative binary operations.
The final declaration \texttt{source\_sign\_counterexample} combines every
hypothesis, odd arity, nonreducibility and failure of the representation.

The complete proof is checked with Lean 4.34.1 and the pinned mathlib
revisions supplied with the certificate. Every local declaration's type
and transitive axioms are audited. The paper, provenance hashes,
independent review and isolated reproduction instructions accompany the
formal source. Only one source question is counted.
\begin{thebibliography}{9}
\bibitem{KS} G. Kiss and G. Somlai, \emph{Associative idempotent nondecreasing
functions are reducible}, arXiv:1707.04341v3, September 4, 2018, Example 4.10
and Conjecture 5.3, equation (24).
\url{https://arxiv.org/abs/1707.04341v3}.
\bibitem{MathDB} MathDB, \emph{Characterization of nonreducible associative
idempotent monotone functions}, problem 336624, accessed October 26, 2026.
\url{https://mathdb.com/p/336624}.
\end{thebibliography}
\end{document}
